% RLJ main.tex Version 2026.1
\documentclass[10pt]{article} % For LaTeX2e

\usepackage{rlj}           % Should be uncommented for submission
%\usepackage[accepted]{rlj} % Should be uncommented for the camera-ready
%\usepackage[preprint]{rlj} % Should be uncommented for preprint versions

%%%%%%%%%%%%%%%%%%%%%%%%%%%%%%%%%%%%%%%%%%%%%%%%%%%%%%%%%%%%%%%%
% Recommended (but not required) packages
%%%%%%%%%%%%%%%%%%%%%%%%%%%%%%%%%%%%%%%%%%%%%%%%%%%%%%%%%%%%%%%%
\usepackage{amssymb}
\usepackage{mathtools}
\usepackage{amsthm}
\usepackage{mathrsfs}
\usepackage{graphicx}
\usepackage{booktabs}
\usepackage{url}

%%%%%%%%%%%%%%%%%%%%%%%%%%%%%%%%%%%%%%%%%%%%%%%%%%%%%%%%%%%%%%%%
% Theorem environments and macros
%%%%%%%%%%%%%%%%%%%%%%%%%%%%%%%%%%%%%%%%%%%%%%%%%%%%%%%%%%%%%%%%
\theoremstyle{definition}
\newtheorem{definition}{Definition}
\theoremstyle{plain}
\newtheorem{theorem}{Theorem}
\newtheorem{lemma}{Lemma}
\newtheorem{proposition}{Proposition}

\newcommand{\Ex}{\mathbb{E}}
\newcommand{\Hist}{\mathcal{H}}
\newcommand{\Acts}{\mathcal{A}}
\newcommand{\ModE}{\mathcal{M}_E}
\newcommand{\PiK}{\Pi_{K_L}}
\newcommand{\Vlow}{V^{\mathrm{lo}}_t}
\newcommand{\Vhi}{V^{\mathrm{hi}}_t}
\newcommand{\Rex}{\mathcal{R}^{\mathrm{ex}}}
\newcommand{\Rmax}{R_{\max}}
\newcommand{\Dstruct}{\Delta^{\mathrm{str}}}

%%%%%%%%%%%%%%%%%%%%%%%%%%%%%%%%%%%%%%%%%%%%%%%%%%%%%%%%%%%%%%%%
% Meta-data (cover page) -- fill in the TODO fields
%%%%%%%%%%%%%%%%%%%%%%%%%%%%%%%%%%%%%%%%%%%%%%%%%%%%%%%%%%%%%%%%
\title{TODO}

\setrunningtitle{TODO}

\author{Anonymous Authors\textsuperscript{1}}

\emails{anon@example.com}

\affiliations{
$^{1}$\textbf{Anonymous Institution}
}

\contribution{TODO}{None}

\keywords{TODO}

\summary{TODO}

\begin{document}

\makeCover
\maketitle

\begin{abstract}
TODO
\end{abstract}

\section{Introduction}
\label{sec:intro}

Reinforcement learning carries an unexamined assumption about evaluation: that an agent is graded \emph{from outside}, against an optimum $V^{*}$ supplied by a model, a prior, or a fixed comparator \citep{lai1985asymptotically,jaksch2010ucrl2}. This is harmless when the world is small enough to model, but continual reinforcement learning (CRL) rests on the opposite premise \citep{ring1994continual,abel2023definition}. Under the \emph{big-world hypothesis} the agent is far smaller than its world, cannot represent or retain a fixed optimal solution, and so never stops learning \citep{javed2025bigworld,lewandowski2025bigger,dohare2024plasticity}. The assumption then fails twice over: at run time the agent has no outside source for $V^{*}$, and even a known $V^{*}$ is the wrong target, since a bounded agent should be judged against the best use of its own resources \citep{simon1955behavioral,russell1995bounded,lewis2014computational,lieder2020resource}. This is the agent-centric turn that recent foundational work urges, judging adaptation rather than the recovery of a fixed solution \citep{abel2024dogmas}.

Concretely, picture an agent on an open-ended stream where new options keep appearing, the situation continual RL is meant to capture. Standard regret asks for the value $V^{*}$ of the best policy; but when the set of behaviours is itself growing, that comparator is undefined at any finite time, and standard regret returns no number. The agent can still ask what it \emph{can} answer from its own stream: how much better do I believe I could be doing than I am? Examiner regret is that question made precise.

We drop the external yardstick and adopt the agent's-eye view: if a continual agent is to know whether it is still learning, it must compute that judgment itself, from its own stream, within its own capacity. We make the judgment an explicit object. Alongside the acting \emph{learner}, an internal \emph{examiner} forecasts an \emph{aspiration} for the best value it believes attainable and a \emph{guarantee} for the value it is sure it achieves; their gap, the \emph{examiner regret} $\rho_t$, is the agent's self-assessed room to improve. Our thesis is that this self-evaluation is the only honest internal measure of progress a big-world agent has, and, the load-bearing claim, that to be trustworthy it must be produced by a module that \emph{out-learns the policy it grades}. Acting needs only a policy, but judging whether one \emph{could} have done better means reasoning about choices never tried, which needs a richer, faster-learning predictive model than the one used to act. Evaluation, in a big world, is not a fixed reference held outside the agent; it is the part of the agent that must keep learning hardest. The price is real, and we state it as a principle of \emph{no free self-evaluation}: an honest self-grade costs either exhaustive exploration or an inductive generalization the agent cannot verify. That price is the classical explore-or-assume dichotomy \citep{kearns2002e3,brafman2002rmax}, so our contribution is not the dichotomy but what it forces on the design of the examiner.

The remainder of the paper develops this position. We define examiner regret and observe that, on its own, it certifies nothing, since an agent can simply declare it to be zero (Proposition~\ref{prop:nothing}). We isolate the missing ingredient as a \emph{soundness} condition and show it makes $\rho_t$ a certificate of the regret a bounded agent can actually reduce (Lemma~\ref{lem:bridge}), recovering the classical bandit rate when the agent can cover its world. We then prove the price: a sound examiner with a vanishing self-report requires either complete coverage of the world or a correct inductive generalization, and neither is available for free in a big world (Theorem~\ref{thm:nofree}). We give the constructive core next, a recipe for making the examiner out-learn the learner on the same stream through off-policy prediction of policies never executed, an asymmetric critic, and a calibration target, followed by an explicit account of the limitations. We claim no new algorithm or rate; the contribution is the position, a definition that makes it precise, a design principle for the examiner, and a lower bound that gives it teeth.

\section{A Learner That Acts and an Examiner That Out-Learns It}
\label{sec:defs}

We work over a single, never-resetting stream of experience and assume no Markov, episodic, or stationarity structure. Rewards lie in $[0,\Rmax]$, and we measure value over an $H$-step look-ahead. The agent is capacity-bounded: it represents policies and beliefs within a fixed budget of bits. We seek a definition of a learning measure the agent can compute from its own state, the condition under which it is trustworthy, and the limit beyond which it cannot be.

\begin{definition}[Stream, policy, value]
\label{def:basic}
A \emph{history} is $h_t = (o_0,r_0,a_1,\dots,a_t,o_t,r_t)$, with observations $o_s$, rewards $r_s \in [0,\Rmax]$, actions $a_s$; the underlying law is $\mu$. A \emph{policy} $\pi$ maps histories to action distributions, with value $f^{\pi}_t = \Ex_{\mu,\pi}\!\big[\sum_{s=t+1}^{t+H} r_s \mid h_t\big]$.
\end{definition}

\begin{definition}[What the learner could have done]
\label{def:capcomp}
The learner represents policies within $K_L$ bits; write $\PiK$ for this representable set. The \emph{within-capacity optimum} is $f^{K_L}_t = \sup_{\pi \in \PiK} f^{\pi}_t$, and the \emph{structural gap} to the global optimum is $\Dstruct_t = f^{*}_t - f^{K_L}_t \ge 0$. We compare the agent to $f^{K_L}_t$, the best it could itself express, and keep $\Dstruct_t$ visible as the irreducible price of being small.
\end{definition}

\begin{definition}[Learner and examiner]
\label{def:agent}
An agent is a pair $(\mathcal{L},\mathcal{E})$. The \emph{learner} $\mathcal{L}$ emits the policy $\pi_t \in \PiK$ it acts with, and learns only about that policy. The \emph{examiner} $\mathcal{E}$ is a \emph{predictive model}: from the history it outputs an \emph{aspiration} $\Vhi$ and a \emph{guarantee} $\Vlow$, with $\Vhi \ge \Vlow$. Setting $\Vhi$ requires forecasting the value of policies the learner did not run, so the examiner must learn about a whole family of behaviours, not the one acted. In the spirit of general value functions \citep{sutton2011horde,modayil2014nexting}, temporally abstract predictions \citep{suttonprecupsingh1999}, and off-policy prediction \citep{precup2001offpolicy}, $\Vhi$ and $\Vlow$ are value predictions the agent learns from its own stream, and the examiner's \emph{inductive bias}, the model class $\ModE$ of environments it considers possible, the concentration rule, or the prior of its predictors, is what lets it generalize to behaviours it has not tried. A Bayesian posterior over a model class is one instance, not the only one \citep{osband2013psrl}.
\end{definition}

\begin{definition}[Examiner regret]
\label{def:exreg}
$\rho_t = \Vhi - \Vlow \ge 0$, with $\Rex_T = \sum_{t=1}^T \rho_t$. It depends only on the history and the examiner; it never refers to $\mu$ or to any optimum.
\end{definition}

\section{A Self-Report Carries No Guarantee}
\label{sec:nothing}

The definition is cheap on purpose, and the first thing to say is what that cheapness costs.

\begin{proposition}[Examiner regret says nothing on its own]
\label{prop:nothing}
There is an examiner that reports $\rho_t \equiv 0$ for every learner and every world. Hence two agents with the identical self-report $\rho_t \equiv 0$ can differ arbitrarily in actual performance: one optimal within its capacity, the other suboptimal by a constant at every step.
\end{proposition}

\begin{proof}
Let the examiner always output $\Vhi = \Vlow = 0$, so $\rho_t \equiv 0$ regardless of behaviour or environment. Pair it once with a learner that plays a within-capacity optimum, and once with a learner that plays a policy worth a constant less at each step. Both report $\rho_t \equiv 0$. \textbf{[proved]}
\end{proof}

So examiner regret guarantees nothing by itself; trust must come from a condition tying the examiner's two numbers to real value. That condition is the subject of the paper, and it is where the examiner's learning enters.

\section{Soundness, and Why the Examiner Must Out-Learn the Learner}
\label{sec:sound}

\begin{definition}[Soundness]
\label{def:sound}
The examiner is \emph{sound} with slack $(\delta_t)$ if, with high probability, for all $t$: the guarantee is a true floor, $\Vlow \le f^{\pi_t}_t$, and the aspiration is a true ceiling for the within-capacity optimum, $\Vhi \ge f^{K_L}_t - \delta_t$. It is \emph{consistent} if the average slack vanishes, $\frac1T\sum_{t\le T}\delta_t \to 0$.
\end{definition}

Soundness rules out the empty report of Proposition~\ref{prop:nothing}: the all-zero examiner violates the ceiling whenever the agent could do better than zero. The substantive question is how an examiner can \emph{meet} the ceiling, since it must set $\Vhi$ above the value of policies the learner may never have run. It has two ways: \emph{coverage}, having seen those policies tried, or \emph{extrapolation} from its inductive bias $\ModE$. This is the precise sense in which the examiner must out-learn the learner. The learner extracts one policy's value from the stream, whereas a sound examiner must extract the value of an entire family, including behaviours never executed. That is strictly more learning from the same data, and it is what a single policy representation cannot supply. The inductive bias is not a liability to hide; it is the ordinary generalization any learner needs, and here it lets a small examiner judge a large space cheaply, at the risk, made precise in Section~\ref{sec:nofree}, of being confidently wrong.

\begin{lemma}[Soundness turns the self-report into a certificate]
\label{lem:bridge}
If the examiner is sound with slack $(\delta_t)$, then with high probability the agent's regret against its within-capacity optimum obeys
\[
\sum_{t=1}^T \big(f^{K_L}_t - f^{\pi_t}_t\big) \;\le\; \Rex_T + \sum_{t=1}^T \delta_t,
\]
and the regret against the global optimum adds only the irreducible $\sum_t \Dstruct_t$.
\end{lemma}

\begin{proof}
For each $t$, write $f^{K_L}_t - f^{\pi_t}_t = (f^{K_L}_t - \Vhi) + (\Vhi - \Vlow) + (\Vlow - f^{\pi_t}_t)$. By soundness the first term is at most $\delta_t$ and the third is at most $0$, while the middle term is $\rho_t$; summing gives the first inequality. Adding $\sum_t \Dstruct_t = \sum_t(f^{*}_t - f^{K_L}_t)$ gives the global statement. \textbf{[proved]}
\end{proof}

This is a decomposition, not a new rate: the identity is the value gap of the performance-difference and optimism arguments familiar from policy improvement and confidence-bound analyses \citep{kakade2002cpi,jaksch2010ucrl2}. \textbf{[cited]} Its point is interpretive. Under soundness the cheap self-report becomes an upper bound on the regret a bounded agent can drive down, with the unremovable structural cost shown rather than hidden.

The certificate is tight when the agent can cover its world. In a $K$-armed bandit whose arms the learner can all represent and try, the examiner needs no extrapolation: the aspiration is the largest upper confidence bound, the guarantee the lower confidence bound of the played arm. By Hoeffding's inequality these are sound with vanishing slack, and the standard UCB1 analysis gives $\Rex_T = O(\sqrt{KT\log T})$ \citep{auer2002ucb}, so by Lemma~\ref{lem:bridge} the within-capacity regret matches the classical rate. \textbf{[cited]} Coverage, not cleverness, makes the slack vanish here, and coverage is exactly what a big world denies.

\section{The Price: No Free Self-Evaluation in a Big World}
\label{sec:nofree}

We now state the price of the position. When the world is larger than the agent can cover, a sound examiner can keep its self-report small only by trusting an inductive generalization, and that generalization cannot be verified, so it can be wrong. Either way the agent pays.

\begin{theorem}[No free self-evaluation]
\label{thm:nofree}
Consider a family of $N$-armed bandits in which exactly one ``special'' arm pays $\tfrac12+\epsilon$ and the rest pay $\tfrac12$; which arm is special is unknown. Run any agent for a horizon $T$. Then on a worst-case member of the family, with high probability,
\[
\rho_t \;\ge\; \epsilon - \delta_t \quad\text{for all } t, \qquad\text{so}\qquad \Rex_T \;\ge\; \epsilon T - \sum_{t\le T}\delta_t,
\]
unless the examiner removes both of the following gaps:
\begin{itemize}
\item[(i)] \emph{a coverage gap}: if $T < N$ the learner leaves some arm untried, and the special arm can be placed there;
\item[(ii)] \emph{a bias gap}: if the examiner's model class $\ModE$ cannot tell the special arm apart from an ordinary one, observing the others does not reveal it.
\end{itemize}
Removing (i) requires trying every arm ($T \ge N$); removing (ii) requires $\ModE$ to contain the truth, which the agent cannot check against arms it has not resolved. Consequently a consistent examiner ($\delta_t \to 0$) that reports $\rho_t \to 0$ in a big world must be relying on an unverified generalization, and if that generalization is wrong it is not sound.
\end{theorem}

\begin{proof}
\emph{Coverage gap.} Let $\mu_0$ be the all-$\tfrac12$ bandit and run the agent in it. In $T$ steps it pulls at most $T$ distinct arms, so if $T < N$ some arm $i^{*}$ is never pulled. Let $\mu_{i^{*}}$ make arm $i^{*}$ special. Since the two worlds differ only at $i^{*}$, which is never pulled, the histories coincide and the examiner outputs the same numbers in both. Under $\mu_{i^{*}}$ the played policy is worth $\tfrac12$, so soundness forces the guarantee $\Vlow \le \tfrac12$ and the aspiration $\Vhi \ge (\tfrac12+\epsilon) - \delta_t$; hence $\rho_t \ge \epsilon - \delta_t$, and summing gives the bound. The agent also never finds the special arm, so its true regret is $\epsilon T$.

\emph{Bias gap.} Suppose instead the learner tries every arm, but the examiner's class $\ModE$ cannot separate the special-arm world from an ordinary one: both are assigned the same value prediction by $\ModE$ (any $\ModE$ that maps the two to a common forecast, for instance a bias asserting the arms share one mean). An examiner that trusts such an $\ModE$ then forecasts the special arm at $\tfrac12$, so it sets $\Vhi \approx \tfrac12 < (\tfrac12+\epsilon) - \delta_t$ and is \emph{unsound} on the world where that arm is special. To restore soundness it must widen the aspiration to cover the possibility it cannot rule out, returning it to a hedge of size $\epsilon$, i.e.\ $\rho_t \ge \epsilon - \delta_t$ again.

In both cases a sound examiner cannot also have $\rho_t < \epsilon - \delta_t$; escaping requires closing the coverage gap (try everything) or the bias gap (a correct, and here unverifiable, model class). \textbf{[proved]}
\end{proof}

The reading is the position made precise. Self-evaluation in a big world is bought with one of two coins: \emph{exploration}, which buys coverage but is impossible when the world outruns the horizon, or \emph{prior knowledge}, which buys extrapolation only if it happens to be right and which the agent cannot audit on what it has not seen \citep{osband2013psrl,russo2017ids}. The examiner's power is exactly the generalization it carries, and it is double-edged: a confident wrong bias produces a small but false self-report. A near-zero $\rho_t$ in a genuinely big world should therefore be read as unearned confidence, not completed learning, which also makes the often-questioned consistency assumption honest, since in a big world it cannot hold cheaply, and the theorem says why.

Both gaps are about \emph{information}, not effort: the coverage gap is missing data, the bias gap a missing distinction in $\ModE$. Making the examiner compute harder over the same stream, by planning, replay, or a larger offline network, can tighten a bound \emph{within} what the data and bias already determine, but cannot manufacture information about an arm never pulled or a distinction $\ModE$ cannot draw. This is why we locate the examiner's power in its learning and data access, not raw computation: a computationally lavish examiner is still bound by Theorem~\ref{thm:nofree}.

The impossibility direction is, deliberately, not novel. Gaps (i) and (ii) are the classical explore-or-assume dichotomy and the realizability requirement, the same forces priced by polynomial-time exploration results \citep{kearns2002e3,brafman2002rmax}, minimax regret lower bounds \citep{azar2017minimax}, and the two-point construction of \citet{lai1985asymptotically}. We invoke it not as a new hardness result but to price the position and constrain the examiner's design, saying exactly what an honest self-grade must buy and hence what the examiner must be built to learn (Section~\ref{sec:practical}).

The favourable bandit and the worst-case family are two sides of one coin: how much of the world the agent can resolve, by exploring or by correctly assuming. The total budget $K = K_L + K_E$ governs both, since bits given to the examiner buy a richer, more reliable predictive model while bits given to the learner buy reach and coverage. Characterizing the split that best trades exploration against generalization, under genuine non-stationarity, is the central problem the position opens.

\section{Making the Examiner Out-Learn the Learner}
\label{sec:practical}

The construction is directly estimable from methods already in use, which makes the position a tool rather than only a stance. \textbf{[heuristic]} The crux a practitioner will press is how the examiner can be \emph{better informed} than the learner when both see the same stream; we answer it concretely here.

The aspiration $\Vhi$ is any \emph{optimistic} value estimate and the guarantee $\Vlow$ any \emph{pessimistic} one, both computed from the same stream, and existing algorithms supply both ends directly (Table~\ref{tab:sources}). In a big-world bandit the natural choice is exactly UCB: $\Vhi$ is the largest upper-confidence index and $\Vlow$ the lower-confidence bound of the played arm \citep{auer2002ucb}, so $\rho_t$ is the confidence width and the examiner's bias is the concentration assumption. In tabular problems they are the optimistic value of confidence-set value iteration \citep{jaksch2010ucrl2,strehl2009pacmdp,azar2017minimax} and a conservative empirical value \citep{kumar2020cql,jin2021pessimism}. In deep RL they are the upper and lower quantiles of a critic ensemble or a bootstrapped value function and a conservative critic \citep{osband2016bootstrapped,lu2021bitbybit,kumar2020cql}. In the Bayesian view they are a posterior upper quantile of the within-capacity optimum and the posterior value of the current policy \citep{osband2013psrl,guez2012bamcp,lidayan2025bampf,arumugam2024capacity}. The gap $\rho_t$ between two such estimates is a close relative of model-value inconsistency as a signal of epistemic uncertainty \citep{filos2022modelvalue}; our addition is the soundness target that turns the signal into a certificate.

\begin{table}[t]
\caption{The examiner's two values are read off existing optimistic and pessimistic estimators; their gap is $\rho_t$, and the estimator's assumption is the examiner's bias $\ModE$.}
\label{tab:sources}
\begin{center}
\resizebox{\textwidth}{!}{%
\small
\begin{tabular}{@{}llll@{}}
\toprule
Setting & Aspiration $\Vhi$ & Guarantee $\Vlow$ & Examiner bias $\ModE$ \\
\midrule
Bandit & UCB index \citep{auer2002ucb} & LCB of played arm & concentration \\
Tabular MDP & optimistic VI \citep{jaksch2010ucrl2} & conservative empirical \citep{jin2021pessimism} & confidence sets \\
Deep RL & ensemble upper \citep{osband2016bootstrapped} & conservative critic \citep{kumar2020cql} & ensemble prior \\
Bayesian & posterior upper quantile \citep{osband2013psrl} & posterior value of $\pi_t$ & model prior \citep{arumugam2024capacity} \\
\bottomrule
\end{tabular}%
}
\end{center}
\end{table}

The learner and examiner see the \emph{same} stream, yet the examiner is strictly more informed along three concrete axes, none of which is ``just have a better prior.'' The first is off-policy counterfactual prediction. The learner learns the value of the one policy it runs, whereas the examiner learns the value of a \emph{family} of policies it does not run, using off-policy temporal-difference learning and general value functions \citep{sutton2011horde,precup2001offpolicy,modayil2014nexting}. Predicting the value of untried behaviour is exactly what setting $\Vhi$ demands, and it is exactly what general value functions and recognizers were built to do; this is the rigorous, non-circular sense in which the examiner extracts more from the same data. The second is asymmetric capacity and timescale. The examiner may carry more capacity and even privileged signals than the real-time learner, as in asymmetric actor--critic \citep{pinto2018asymmetric}, and may be driven to track on a faster timescale than the policy it evaluates, as in two-timescale actor--critic and stochastic approximation \citep{konda2000actorcritic,borkar1997twotimescale}, so that the evaluation stays ahead of the behaviour. The third is a calibration target. Soundness is the requirement that the examiner's estimates be valid on what has been seen, so training the examiner to be calibrated on covered behaviour \citep{foster1998calibration} makes ``is the examiner learning?'' operational, since its slack $\delta_t$ shrinks on the covered region. The reading rule follows: a \emph{sound, shrinking} $\rho_t$ is certified progress, a \emph{small but unsound} $\rho_t$ is unearned confidence, and Theorem~\ref{thm:nofree} says when the latter is unavoidable. None of the three manufactures information, so the examiner's lead is real but capped by coverage, the content of the lower bound, now a design rule: separate the critic from the actor and give it its own larger budget, off-policy targets, and a calibration loss.

The construction is then directly usable in three ways. As an \emph{exploration signal}, $\rho_t$ fed back as an intrinsic reward drives the agent toward states where optimism and pessimism disagree, closing the otherwise-absent loop between the modules; it is a value-unit relative of learning-progress and prediction-error drives \citep{schmidhuber2010formal,oudeyer2007intrinsic,bellemare2016count,pathak2017curiosity,burda2019rnd}. As a \emph{comparison metric}, comparing the $\rho_t$ trajectories of different algorithms on one continual stream estimates which is nearer its own capacity-constrained optimum without any $V^{*}$. And since soundness is calibration on what has been seen, a practitioner can \emph{audit} a low $\rho_t$, while Theorem~\ref{thm:nofree} says precisely when no audit is possible.

A concrete deployment makes the value plain. Consider monitoring a deployed continual agent, a recommender or a robot, that one cannot grade against an optimum. The operator wants to know whether it has genuinely stopped improving or is merely stuck. Examiner regret gives an online read of self-assessed room to improve, auditable on covered behaviour, and Theorem~\ref{thm:nofree} is the warning the operator needs: a confidently small $\rho_t$ may be unearned, so apparent convergence should trigger exploration or external testing rather than trust. We do not run these experiments; we map the path because it is what turns the position into practice.

\section{Limitations}
\label{sec:limits}

\textbf{(L1)} The content lives in soundness, not in $\rho_t$: Proposition~\ref{prop:nothing} is explicit that the self-report alone certifies nothing, and soundness is extrinsic, requiring the examiner to carry a model class and hence a prior. \textbf{(L2)} Lemma~\ref{lem:bridge} is a decomposition, the optimism gap relabeled, not a new rate; its value is conceptual. \textbf{(L3)} Our tight, positive result is shown only for the coverable bandit; soundness of a biased, capacity-bounded examiner under genuine non-stationarity, and the under-capacity case $\Dstruct_t > 0$, are not established here. \textbf{[GAP]} \textbf{(L4)} Theorem~\ref{thm:nofree} is a worst-case statement over a family, in the manner of two-point lower bounds \citep{lai1985asymptotically,azar2017minimax}; it does not bound a fixed environment, and as we note in Section~\ref{sec:nofree} the impossibility direction is not itself novel. The bias gap holds for any $\ModE$ that cannot separate the two worlds, with exchangeability as one instance, but we have not characterized which $\ModE$ avoid it on which world families. \textbf{(L5)} The out-learning claim of Section~\ref{sec:practical} is a design argument supported by existing off-policy and asymmetric-critic methods, not a theorem: we do not prove that an off-policy examiner attains soundness faster than the learner attains optimality in general. \textbf{[GAP]} \textbf{(L6)} In the bare formalism the learner and examiner are separate maps; we propose one coupling, feeding $\rho_t$ back as an intrinsic reward, but we do not analyze the resulting joint dynamics, and the optimal budget split $K = K_L + K_E$ remains conceptual. \textbf{(L7)} The capacity model counts representable policies and treats compute only informally; a joint account of description length and computation is left open. \textbf{(L8)} Relative to capacity-limited Bayesian RL \citep{arumugam2024capacity}, which already uses within-capacity targets, and to model-value inconsistency \citep{filos2022modelvalue}, which already reads a value-estimate gap as uncertainty, the novelty is the explicit examiner, the self-report-certifies-nothing thesis, the out-learning design principle, and the coverage-or-bias lower bound, not a new upper bound. \textbf{(L9)} There are no experiments; the recipe of Section~\ref{sec:practical} is untested.

\section{Related Work}
\label{sec:related}

\textbf{Foundations and bounded rationality.} \citet{ring1994continual,abel2023definition,abel2024dogmas,abel2020bounded} define and develop CRL and the agent-centric, adaptation-first view, and \citet{dohare2024plasticity} document how deep learners lose the very plasticity continual learning needs. Bounded optimality \citep{simon1955behavioral,russell1995bounded}, computational rationality \citep{lewis2014computational}, and resource-rational analysis \citep{lieder2020resource} motivate judging a bounded agent against its own resources, our comparator $f^{K_L}$; the reward hypothesis \citep{silver2021reward} frames the goal that examiner regret measures progress toward.
\textbf{Big-world hypothesis.} \citet{javed2025bigworld,lewandowski2025bigger} argue the agent is far smaller than its world, the premise behind our removal of the external yardstick.
\textbf{Predictive knowledge as the examiner's substrate.} General value functions and Horde \citep{sutton2011horde,modayil2014nexting}, temporal abstraction and options \citep{suttonprecupsingh1999}, off-policy temporal-difference learning \citep{precup2001offpolicy}, asymmetric actor--critic \citep{pinto2018asymmetric}, and two-timescale learning \citep{konda2000actorcritic,borkar1997twotimescale} are the machinery by which an evaluation module learns about behaviours it does not enact; we use them to make ``out-learn the learner'' concrete.
\textbf{Optimism, pessimism, and exploration.} Confidence-bound and PAC analyses \citep{kearns2002e3,brafman2002rmax,jaksch2010ucrl2,auer2002ucb,strehl2009pacmdp,azar2017minimax} and the performance-difference lemma \citep{kakade2002cpi} supply the optimistic aspiration; conservative and pessimistic value learning \citep{kumar2020cql,jin2021pessimism} and bootstrapped value functions \citep{osband2016bootstrapped} supply the guarantee.
\textbf{Bayesian and information-theoretic RL.} \citet{osband2013psrl,russo2017ids} encode the prior the examiner extrapolates from; BAMDP and AIXI \citep{guez2012bamcp,lidayan2025bampf,hutter2005aixi} supply the model-class examiner and its idealized limit; \citet{lu2021bitbybit,arumugam2024capacity} cast capacity in bits, the latter being the closest prior work.
\textbf{Learning progress, self-assessment, and certification.} Self-assessed and prediction-error drives \citep{schmidhuber2010formal,oudeyer2007intrinsic,bellemare2016count,pathak2017curiosity,burda2019rnd} and model-value inconsistency \citep{filos2022modelvalue} read internal disagreement as a signal; calibration \citep{foster1998calibration}, KWIK's ``knows what it knows'' \citep{li2011kwik}, non-stationary regret \citep{wei2021master}, and hindsight rationality \citep{morrill2021hindsight} bear on honest self-assessment. Across all of these the evaluator is external or fused with the policy, whereas we make it an explicit, continually-learning, capacity-bounded module and ask what its self-report is worth.

\section{Discussion}
\label{sec:discussion}

Continual reinforcement learning has no examiner standing outside the agent, so the regret that matters is the one a bounded internal examiner can certify. We have argued, and made precise, that such self-evaluation is the only honest internal measure of progress, is meaningful only under soundness, and that meeting soundness forces the examiner to be the module that out-learns the policy, because judging counterfactual behaviour means learning about choices never tried, which a single acting policy cannot supply. In a big world this is never free: a trustworthy self-report demands either coverage the horizon disallows or a generalization the agent cannot verify. The picture turns evaluation in CRL into a budgeting problem between acting and judging, makes ``is the agent still learning?'' an operational question answered by a sound, shrinking $\rho_t$, and warns against trusting an agent that reports it has finished learning. We claim no new algorithm or rate; the contribution is the position, a definition, a design principle for the examiner, and a lower bound that shows the price is real.

\bibliography{main}
\bibliographystyle{rlj}

\end{document}
